\begin{abstract}
We introduce \emph{tactile analogies}, an example-based synthesis paradigm for tactile videos.
Given a small set of reference touches recorded on an object, we predict the tactile signal that
a vision-based tactile sensor would measure at a new, previously untouched location. We represent
a touch as a \emph{tactile normal video}, which records how the deformable gel surface of the
sensor bends over the course of a press, and we bridge reference and query locations through
surface normal maps rendered at the two sensor poses over several fields of view. Prediction
proceeds coarse to fine: a self-supervised image descriptor retrieves the most promising
reference touch, a pre-trained local feature matcher aligns the two normal renderings and warps
the reference video into the query frame, and a lightweight video network refines the warped
video into the final prediction. To evaluate the task we build a benchmark on the ObjectFolder
dataset with two settings, one that isolates alignment and refinement and one that exercises the
full pipeline including retrieval. Our method outperforms baselines adapted from prior work in
tactile prediction while running fast enough for interactive use. Finally, we show that the
predicted normals can be integrated into heightmaps for fine-grained 3D surface reconstruction,
and pushed through an optical tactile simulator to produce virtual visuo-tactile measurements at
locations that were never physically touched.
\end{abstract}
